\chapter{RV32I assembly I: arithmetic, memory, control}
\label{ch:03}

Chapter~2 showed that the machine sees only bytes. This chapter shows
what it \emph{does} with them: the 37 instructions of RV32I, the base
RISC-V integer set --- and the instruction set your capstone CPU will
execute for real in chapters 5--8. Everything here is concrete: a small
assembler you can read end-to-end encodes each program to real hex, a
small interpreter runs it, and every number in the chapter comes from
that run --- not from a datasheet summary. If you can write a loop in
this chapter's notation, you can read the disassembly of any program
your laptop runs.

\index{RV32I}\index{assembly language}\index{register file}%
\index{ABI names}\index{instruction encoding}\index{R-type}%
\index{I-type}\index{S-type}\index{B-type}\index{U-type}\index{J-type}%
\index{pseudoinstruction}\index{branch}\index{load}\index{store}

\section{The register file and the ABI names}
\label{sec:regs03}

\index{x0}\index{zero register}
Thirty-two registers, \texttt{x0}--\texttt{x31}, each 32 bits --- the
same file shape as your cpu8's, just wider and deeper. Two of them are
special by hardware: \texttt{x0} is hardwired to zero (writes to it are
silently discarded, which is why \texttt{mv rd, rs} is really
\texttt{addi rd, rs, 0} and \texttt{nop} is really
\texttt{addi x0, x0, 0}). The rest are general-purpose by hardware but
assigned jobs by convention --- the ABI names you will see in every
disassembly:

\begin{figure}[htbp]
\centering
\begin{center}
\begin{tabularx}{0.95\linewidth}{llX}
\toprule
\textbf{ABI name} & \textbf{x-name} & \textbf{Job} \\
\midrule
\texttt{zero} & x0 & wired to 0 --- the register that isn't \\
\texttt{ra} & x1 & return address \\
\texttt{sp} & x2 & stack pointer \\
\texttt{gp}, \texttt{tp} & x3, x4 & globals, thread pointer \\
\texttt{t0}--\texttt{t2} & x5--x7 & temporaries (caller-saved) \\
\texttt{s0}/\texttt{fp}, \texttt{s1} & x8, x9 & saved / frame pointer \\
\texttt{a0}--\texttt{a7} & x10--x17 & args + return values \\
\texttt{s2}--\texttt{s11} & x18--x27 & saved (callee-saved) \\
\texttt{t3}--\texttt{t6} & x28--x31 & more temporaries \\
\bottomrule
\end{tabularx}
\end{center}

\caption{The RV32I register file with ABI conventions. Programs obey
the hardware; \emph{programmers} obey this table so calls compose.}
\label{fig:regs03}
\end{figure}

\begin{principle}
The ABI names are a contract between functions, not hardware.
\texttt{a0}--\texttt{a7} carry arguments and return values;
\texttt{t0}--\texttt{t6} are temporaries a callee may trash;
\texttt{s0}--\texttt{s11} a callee must preserve. ch4 makes the stack
discipline formal; for now: use \texttt{t}-names inside one function
and the world stays simple.
\end{principle}

\section{Six formats, one discipline}
\label{sec:formats03}

\index{opcode}\index{funct3}\index{funct7}\index{immediate}
Every RV32I instruction is 32 bits --- fixed width, no exceptions ---
carved into six layouts. The carving is deliberately regular: register
fields (\texttt{rd}, \texttt{rs1}, \texttt{rs2}), the function selector
\texttt{funct3}, and the 7-bit \texttt{opcode} keep the same bit
positions across every format they appear in. Only the immediate bits
move, and even they preserve their internal order. That regularity is
not elegance for its own sake --- it is what lets a decoder be a few
hundred gates rather than a table lookup.

\begin{figure}[htbp]
\centering
\begin{tikzpicture}[font=\tiny, >=latex]
% six formats: each row = name + six labelled boxes
\newcommand{\fmtrow}[8]{%
  \node[anchor=east, font=\scriptsize\bfseries] at (0.5,#1) {#2};
  \draw[thick, fill=graybg]     (0.6,#1-0.27) rectangle ++(1.55,0.55) node[midway]{#3};
  \draw[thick, fill=accentsoft] (2.15,#1-0.27) rectangle ++(1.55,0.55) node[midway]{#4};
  \draw[thick, fill=accentsoft] (3.7,#1-0.27) rectangle ++(1.55,0.55) node[midway]{#5};
  \draw[thick, fill=graybg]     (5.25,#1-0.27) rectangle ++(1.55,0.55) node[midway]{#6};
  \draw[thick, fill=goodbg]     (6.8,#1-0.27) rectangle ++(1.55,0.55) node[midway]{#7};
  \draw[thick, fill=accentsoft] (8.35,#1-0.27) rectangle ++(1.55,0.55) node[midway]{#8};}
\fmtrow{5.0}{R}{f7}{rs2}{rs1}{f3}{rd}{op}
\fmtrow{4.0}{I}{imm[11:0]}{}{rs1}{f3}{rd}{op}
\fmtrow{3.0}{S}{imm[11:5]}{rs2}{rs1}{f3}{imm[4:0]}{op}
\fmtrow{2.0}{B}{imm[12|10:5]}{rs2}{rs1}{f3}{imm[4:1|11]}{op}
\fmtrow{1.0}{U}{imm[31:12]}{}{}{}{rd}{op}
\fmtrow{0.0}{J}{imm[20|10:1|11|19:12]}{}{}{}{rd}{op}
\node[align=left, text=slate, text width=10cm] at (5.2,-0.8)
  {Same 32 bits, six layouts. Notice the discipline: rd, rs1, rs2, f3
   and opcode sit in the same places in every format --- only the
   immediate fields move, and even they keep their bit order. The
   encoding exists to make the decoder cheap.};
\end{tikzpicture}

\caption{The six encoding formats. Green = destination, blue =
source/opcode, grey = immediate fields.}
\label{fig:formats03}
\end{figure}

\begin{example}[Hand-assemble, then verify against the real assembler]
\label{ex:asm03}
\texttt{addi a0, a0, 1}: I-type, imm=1, rs1=\texttt{a0}=x10,
\texttt{f3}=0, rd=x10, op=\texttt{0x13} ---
$1\ll20 \mid 10\ll15 \mid 10\ll7 \mid \texttt{0x13}
= \texttt{0x00150513}$. The book's assembler
(\texttt{code/ch03\_rv32asm.py}, a 40-line Python RV32I assembler)
produces the same word --- real output on the five-format smoke test:

\begin{lstlisting}[basicstyle=\ttfamily\scriptsize]
00000000: 00150513   # addi a0, a0, 1
00000004: 00b50863   # beq a0, a1, done
00000008: 00b53633   # add a2, a0, a1
0000000c: 00812683   # lw a3, 8(sp)
00000010: 00d12623   # sw a3, 12(sp)
00000014: 00000013   # nop
\end{lstlisting}

Check \texttt{beq} by hand: B-type needs the offset $16$ (bytes, in
2-byte units in the encoding: imm=16). The assembler computed
\texttt{done}-pc = $0x14-0x04 = 16$ --- and the encoding
\texttt{0x00b50863} decomposes to exactly that.
\end{example}

\begin{tripwire}[Pseudoinstructions are lies you tell on purpose]
\texttt{li}, \texttt{mv}, \texttt{nop}, \texttt{j}, \texttt{ret},
\texttt{bgt} --- none exist in hardware. The assembler expands each to
a real instruction (\texttt{li t0,1} $\to$ \texttt{addi t0,x0,1}). That
is friendly until debugging shows \texttt{nop} where you wrote a
pseudoinstruction, or \texttt{li} becomes two instructions for a big
constant. Read the listing, not the source.
\end{tripwire}

\section{Loads, stores and addressing}
\label{sec:ldst03}

\index{base+offset addressing}
RV32I has exactly one addressing mode: \emph{base register $+$
sign-extended 12-bit immediate} --- \texttt{lw rd, off(rs1)} and
\texttt{sw rs2, off(rs1)}. There is no scaled index, no auto-increment,
no memory-to-memory op. Walking an array is two instructions: the load,
and the pointer bump. That austerity is deliberate: one mode keeps the
address adder single, and the ALU does everything else.

\begin{figure}[htbp]
\centering
\begin{tikzpicture}[font=\scriptsize, >=latex]
% load/store addressing
\draw[thick, fill=accentsoft] (0.4,3.2) rectangle (2.0,4.4);
\node[align=center] at (1.2,4.0) {rs1};
\node[align=center, font=\tiny] at (1.2,3.5) {base = 0x100};
\node[font=\Large] at (2.6,3.8) {$+$};
\draw[thick, fill=graybg] (3.0,3.2) rectangle (4.2,4.4);
\node[align=center] at (3.6,4.0) {imm[11:0]};
\node[align=center, font=\tiny] at (3.6,3.5) {offset = 4};
\draw[->, very thick] (4.6,3.8) -- (5.6,3.8);
\draw[thick, fill=goodbg] (6.0,3.2) rectangle (9.2,4.4);
\node[align=center] at (7.6,4.0) {address 0x104};
\node[align=center, font=\tiny] at (7.6,3.5) {reads/writes word a[1]};
\node[align=left, text=slate, text width=10cm] at (5.0,2.3)
  {\ttfamily lw t4, 4(t0) \rmfamily is the whole addressing story:
   register $+$ sign-extended 12-bit immediate. No indexing modes,
   no scaling --- RV32 keeps the ALU cheap and lets the pointer-bump
   instruction do the walking (that's the \ttfamily addi t0,t0,4
   \rmfamily in the loop).};
\end{tikzpicture}

\caption{The single addressing mode: register $+$ 12-bit immediate.}
\label{fig:ldst03}
\end{figure}

\begin{example}[An array sum, real output]
\label{ex:arrsum03}
\texttt{code/ch03\_arrsum.S} stores \{1,2,3,4,5\} at \texttt{0x100},
then walks it: \texttt{lw t4, 0(t0)} / \texttt{addi t0,t0,4} per
element. The interpreter (\texttt{ch03\_rv32sim.py}, a $\sim$60-line
RV32I simulator with a memory-mapped output register at
\texttt{0x10000000} --- the same trick Book~5 ch16's UART used) runs
the assembled words and prints the real result:

\begin{lstlisting}[basicstyle=\ttfamily\scriptsize]
15
# 56 instructions retired
\end{lstlisting}
\end{example}

\begin{checkpoint}
After \texttt{lw t4, 4(t0)} with \texttt{t0=0x100}, which byte
addresses are read, and what changes if you wrote \texttt{lb} instead?
\end{checkpoint}
\begin{solution}
\texttt{lw} reads \texttt{0x104}--\texttt{0x107} (one aligned word).
\texttt{lb} reads only \texttt{0x104} and sign-extends its byte ---
same base, different width, different sign rule.
\end{solution}

\section{Branches and comparisons}
\label{sec:branch03}

\index{conditional branch}
Six conditional branches: \texttt{beq}, \texttt{bne}, \texttt{blt},
\texttt{bge}, \texttt{bltu}, \texttt{bgeu} --- compare two registers,
jump by a signed immediate if true. There is no flags register: the
compare happens inside the branch, and the ``set'' pseudo-ops
(\texttt{slt}, \texttt{sltu}) exist only when you need the boolean as
data. Each branch is one decision: taken $\to$ \texttt{pc $+$ imm},
not taken $\to$ \texttt{pc $+$ 4}.

\begin{figure}[htbp]
\centering
\begin{tikzpicture}[font=\scriptsize, >=latex]
% branch: pc next decision
\draw[thick, fill=accentsoft] (0.4,2.2) rectangle (3.2,3.4);
\node[align=center, font=\tiny] at (1.8,2.9) {\ttfamily bge t2, t0, loop\\at pc = 0x14};
\draw[->, thick, bend right=20] (1.8,2.2) to node[left, font=\tiny]{taken: pc $\leftarrow$ 0x14 + (-8)} (0.8,0.9);
\draw[thick, fill=goodbg] (0.2,0.4) rectangle (1.6,0.9);
\node[font=\tiny] at (0.9,0.65) {pc = 0x0c};
\draw[->, thick, bend left=20] (1.8,2.2) to node[right, font=\tiny]{not taken: pc $\leftarrow$ pc + 4} (3.2,0.9);
\draw[thick, fill=graybg] (2.8,0.4) rectangle (4.2,0.9);
\node[font=\tiny] at (3.5,0.65) {pc = 0x18};
\node[align=left, text=slate, text width=10cm] at (5.2,-0.6)
  {The branch's whole job: pick the next pc. The compare (bge = rs1
   $\geq$ rs2) decides between pc+4 and pc+imm --- one adder, one mux,
   same shape as the FSM branch on your cpu8.};
\end{tikzpicture}

\caption{A branch picks the next pc: $+4$ or $+$imm --- one mux.}
\label{fig:pcnext03}
\end{figure}

\begin{example}[\texttt{max(17, 42)} with \texttt{bge}]
\label{ex:max03}
\texttt{code/ch03\_max.S}: if \texttt{a0 $\geq$ a1} the answer is
\texttt{a0}, else \texttt{a1}. The branch is one instruction; the
interpreter's run:

\begin{lstlisting}[basicstyle=\ttfamily\scriptsize]
42
# 27 instructions retired
\end{lstlisting}
\end{example}

\begin{checkpoint}
In \texttt{bge t2, t0, loop} at \texttt{pc=0x14} with \texttt{t2=10},
\texttt{t0=5}: is the branch taken, and to which pc?
\end{checkpoint}
\begin{solution}
Taken: $10 \geq 5$, so \texttt{pc} becomes $0x14 + (-8) = 0x0c$ ---
the loop header.
\end{solution}

\section{Assembler directives and running it}
\label{sec:run03}

\index{directives}\index{CPUlator}
Directives are instructions to the assembler, not to the machine:
\texttt{.text}/\texttt{.data} pick the section, \texttt{.word}/\texttt{.byte}
plant data, \texttt{.globl} exports a symbol, labels name addresses.
The chapter's programs print through a memory-mapped output register,
which is what a bare-metal UART driver looks like before an OS arrives.
Two ways to run the same bytes: \textbf{CPUlator} (browser, free) draws
the register file and memory live while single-stepping ---
Figure~\ref{fig:cpulator03} maps its panels; and the book's own
\texttt{ch03\_rv32sim.py} runs the identical words headless (all
outputs in this chapter are from it). Spike is the canonical
command-line emulator; either way the bytes are the same.

\begin{figure}[htbp]
\centering
\begin{tikzpicture}[font=\scriptsize, >=latex]
% CPUlator panel map (schematic, not screenshot)
\draw[thick, rounded corners] (0.3,0.6) rectangle (9.9,4.8);
\node[font=\footnotesize\bfseries] at (5.1,4.5) {CPUlator (browser RV32 simulator) --- panel map};
\draw[thick, fill=accentsoft] (0.6,1.4) rectangle (3.0,4.2);
\node[align=center, font=\tiny, text width=2.2cm] at (1.8,3.8) {editor:\\your .S here};
\node[align=center, font=\tiny, text width=2.2cm] at (1.8,2.4) {compile\\errors land here};
\draw[thick, fill=goodbg] (3.3,1.4) rectangle (5.6,4.2);
\node[align=center, font=\tiny, text width=2.1cm] at (4.45,3.8) {register file:\\x0--x31 live};
\node[align=center, font=\tiny, text width=2.1cm] at (4.45,2.4) {watch t0 count\\to 11, t1 to 55};
\draw[thick, fill=graybg] (5.9,1.4) rectangle (8.2,4.2);
\node[align=center, font=\tiny, text width=2.1cm] at (7.05,3.8) {memory view:\\text + data hex};
\node[align=center, font=\tiny, text width=2.1cm] at (7.05,2.4) {find the loop\\bytes you assembled};
\draw[thick, fill=amberbg] (8.5,1.4) rectangle (9.6,4.2);
\node[align=center, font=\tiny, text width=0.9cm] at (9.05,2.8) {step\\run\\reset};
\node[align=left, text=slate, text width=10cm] at (5.1,0.4)
  {Schematic of the CPUlator layout (CPUlator runs in a browser ---
   described, not captured). Single-step the sum loop and watch the
   same registers this chapter uses change one word at a time.};
\end{tikzpicture}

\caption{CPUlator's four panels, schematically. Described, not
captured --- the tool is a browser app.}
\label{fig:cpulator03}
\end{figure}

\subsection{Measurement: dynamic instruction count}
\label{sec:measure03}

The same sum-1..10 loop written two ways: register variables vs the
``$-O0$ habit'' of keeping \texttt{i} and \texttt{s} in memory. Real
counts from \texttt{ch03\_rv32sim.py}:

\begin{lstlisting}[basicstyle=\ttfamily\scriptsize]
memory-bound (i, s in RAM):     110 instructions retired
register version (i, s in t0/t1): 57 instructions retired
\end{lstlisting}

\begin{figure}[htbp]
\centering
\begin{tikzpicture}[font=\scriptsize, >=latex]
% dynamic instruction count: mem-bound vs register version (real)
\node[align=left, font=\footnotesize\bfseries] at (4.5,4.6)
  {instructions retired, sum-1..10 (real sim runs)};
\draw[thick, fill=accentsoft] (1.2,1.0) rectangle (3.4,4.0);
\node[font=\footnotesize] at (2.3,2.5) {110};
\node[align=center, font=\tiny] at (2.3,0.7) {memory-bound\\(i, s in RAM)};
\draw[thick, fill=goodbg] (5.6,1.0) rectangle (7.8,3.0);
\node[font=\footnotesize] at (6.7,2.0) {57};
\node[align=center, font=\tiny] at (6.7,0.7) {register version\\(i, s in t0, t1)};
\draw[->, thick, slate] (1.2,1.0) -- (8.6,1.0);
\node[align=left, text=slate, text width=10cm] at (4.8,-0.4)
  {Predicted $\sim$2$\times$, measured 1.93$\times$. The loop body is
   7 instructions vs 3 --- keeping variables in registers is not an
   optimisation, it is the difference between code and bookkeeping.
   Counted by \ttfamily ch03\_rv32sim.py\rmfamily, not estimated.};
\end{tikzpicture}

\caption{Real dynamic counts: registers vs memory for the same loop.}
\label{fig:insn03}
\end{figure}

Predicted $\sim$2$\times$ before running (3-instr loop body vs
7-instr); measured $110/57 = 1.93\times$. The gap is the loads and
stores themselves: memory-keeping adds \texttt{lw}/\texttt{sw} traffic
that computes nothing --- bookkeeping billed as instructions. That is
the -O0/-O2 difference you will meet in ch4, visible here without a
compiler.

\begin{wins}
\begin{itemize}
\item You can read any RV32I line and name its format, its fields and
its effect --- and write the 32-bit word for a handful of them.
\item You know the register file's two contracts: hardware
(\texttt{x0}=0) and ABI (\texttt{a}-args, \texttt{t}-temps,
\texttt{s}-saved).
\item You can trace a loop's pc through a branch and predict the next
fetch.
\item You can run RV32I here (book simulator) or in the browser
(CPUlator) and count what it costs.
\item You can point at an instruction stream and say which 4 bytes the
CPU will fetch next, and why.
\end{itemize}
\end{wins}

\section{Common mistakes}
\label{sec:mistakes03}

\begin{itemize}
  \item \textbf{Counting pseudoinstructions.} \texttt{li}/\texttt{mv}/
        \texttt{nop} are sugar: what retires is the expansion. Dynamic
        counts only see real instructions.
  \item \textbf{Forgetting \texttt{x0} eats writes.}
        \texttt{addi x0, x5, 0} does nothing; the register file is
        honest that way.
  \item \textbf{Immediates are 12 bits.} \texttt{addi t0, t0, 4096}
        won't assemble --- big constants need \texttt{lui}$+$\texttt{addi}
        or \texttt{li} (which expands).
  \item \textbf{Branches move the pc, not flags.} There is no carry
        flag to ``check'' --- each branch recomputes its compare.
  \item \textbf{\texttt{lw} where \texttt{lb} belonged.} Byte vs word
        is a correctness issue on packed data, not a performance one.
  \item \textbf{``Memory ops are free.''} The 110-vs-57 measurement is
        the counterexample: the tax is the instruction count itself.
\end{itemize}

\section*{Exercises}
\label{sec:exercises03}
\addcontentsline{toc}{section}{Exercises}

\begin{exercise}\label{ex:e03a}
\texttt{addi x0, x5, 0} --- what does this write, and why is it a
useful instruction anyway?
\end{exercise}
\begin{solution}[hint]
Nothing: \texttt{x0} discards writes. It's \texttt{nop} --- a 4-byte
``do nothing'' that real pipelines need.
\end{solution}

\begin{exercise}\label{ex:e03b}
\texttt{lw a0, 8(sp)} decodes to \texttt{0x00812503}? Verify or
correct, field by field.
\end{exercise}
\begin{solution}[hint]
I-type: imm=8$\ll$20=\texttt{0x00800000}; rs1=\texttt{sp}=x2$\ll$15;
f3=2; rd=\texttt{a0}=x10$\ll$7; op=\texttt{0x03}. The shown word is
right.
\end{solution}

\begin{exercise}\label{ex:e03c}
In the register/ABI table: which two register groups may a callee
freely destroy, and which must it restore?
\end{exercise}
\begin{solution}[hint]
Caller-saved: \texttt{t0}--\texttt{t6} (and args). Callee-saved:
\texttt{s0}--\texttt{s11}. The \texttt{t}/\texttt{s} split is the
whole discipline.
\end{solution}

\begin{exercise}\label{ex:e03d}
\texttt{beq} targets are limited to about $\pm4$~KB. Why, and what
real instruction escapes the limit?
\end{exercise}
\begin{solution}[hint]
B-type holds a 13-bit signed byte offset. Long jumps use \texttt{jal}
(J-type, $\pm1$~MB) or \texttt{jalr} (register-indirect, anywhere).
\end{solution}

\begin{exercise}\label{ex:m03a}
Hand-decode \texttt{0xfe53dce3} (the loop branch) into fields and
offset.
\end{exercise}
\begin{solution}[hint]
B-type; f3=5 (\texttt{bge}); imm bits give $-8$; rs1=\texttt{t2},
rs2=\texttt{t0}. At \texttt{pc=0x14} the target is \texttt{0x0c}.
\end{solution}

\begin{exercise}\label{ex:m03b}
The array-sum walks words with \texttt{addi t0,t0,4}. Rewrite one
iteration to use the immediate instead of the pointer bump, and name
the trade-off.
\end{exercise}
\begin{solution}[hint]
\texttt{lw t4, 4(t0)}... unrolls offsets but keeps the base fixed ---
fewer instructions per element, more code bytes, and it caps at
imm=2048 words. Classic size-vs-count trade.
\end{solution}

\begin{exercise}\label{ex:m03c}
\texttt{slt a0, a1, a2} is ``set if less-than''. Write the two-line
sequence that implements \texttt{max} \emph{without} any branch.
\end{exercise}
\begin{solution}[hint]
\texttt{slt t0, a0, a1} then select: branchless needs arithmetic on
the flag (\texttt{sub}/mask trick) or a conditional move RV32I
doesn't have --- branchless \texttt{max} on base RV32I costs a
multiply or a bit-trick; name it.
\end{solution}

\begin{exercise}\label{ex:m03d}
The memory-bound loop runs 7 instructions per iteration. Count them
from the listing and identify which two are pure bookkeeping.
\end{exercise}
\begin{solution}[hint]
lw i, lw s, add, sw s, addi i, sw i, bge = 7. The two
\texttt{sw}s (and arguably the \texttt{lw}s) compute nothing ---
they only persist state registers would have held for free.
\end{solution}

\begin{exercise}\label{ex:m03e}
\texttt{li a7, 0x10000000} won't fit \texttt{addi}. Write the
two-instruction reality the assembler expands to.
\end{exercise}
\begin{solution}[hint]
\texttt{lui a7, 0x10000} $+$ \texttt{addi a7, a7, 0} --- U-type
plants the top 20 bits, I-type patches the bottom 12.
\end{solution}

\begin{exercise}[hard]\label{ex:h03a}
Implement \texttt{strlen}, \texttt{memcpy} and \texttt{memset} in
RV32I and self-check on the book simulator: fill a buffer, copy it,
measure it, print the length.
\end{exercise}
\begin{solution}[hint]
All three are the same loop skeleton: \texttt{lb}/\texttt{sb} plus a
pointer bump and a compare. \texttt{strlen} counts until 0;
\texttt{memcpy} copies until $n$; \texttt{memset} fills until $n$.
Watch for the NUL byte in the length count.
\end{solution}

\begin{exercise}[hard]\label{ex:h03b}
Extend \texttt{ch03\_rv32asm.py} with \texttt{jal}/\texttt{jalr}
(the J-type format). Encode \texttt{jal ra, +8} and decode it back.
\end{exercise}
\begin{solution}[hint]
J-type packs the immediate in bit-order pieces like B but shifted;
the adder-tree trick is the same. Verify $+8$ lands in the right
fields.
\end{solution}

\begin{exercise}[hard]\label{ex:h03c}
Write a program that prints the first 10 Fibonacci numbers on the
book simulator's MMIO output. Count retired instructions for the
register version vs a memory version.
\end{exercise}
\begin{solution}[hint]
Fib($n$)=Fib($n{-}1$)+Fib($n{-}2$) needs two live registers and one
temp per step --- the register win is bigger than the sum loop's
because there are no array loads at all.
\end{solution}

\section*{Summary}
\label{sec:summary03}
\addcontentsline{toc}{section}{Summary}

\begin{center}
\begin{tabularx}{\linewidth}{lX}
\toprule
\textbf{Idea} & \textbf{One-line version} \\
\midrule
register file & 32$\times$32; \texttt{x0}=0; ABI names split caller-/
callee-saved \\
six formats & rd/rs/f3/op fixed; only immediate fields shuffle \\
addressing & one mode: reg $+$ 12-bit signed immediate \\
branches & compare+offset in one; pc $+$ 4 or pc $+$ imm \\
pseudoinstructions & sugar; count the expansion, not the line \\
measurement & memory-bound loop: 110 vs 57 retired --- the -O0 tax \\
\bottomrule
\end{tabularx}
\end{center}

\section*{Further reading}
\label{sec:fr03}
\addcontentsline{toc}{section}{Further reading}

\begin{itemize}
  \item RISC-V unprivileged ISA specification --- the authoritative
        encoding tables this chapter compresses.
        \url{https://riscv.org/specifications/}
  \item Hennessy \& Patterson, \emph{Computer Organization and Design},
        RISC-V edition, ch.\,2 --- the classic ISA walk-through.
  \item Bryant \& O'Hallaron, \emph{CS:APP}, ch.\,3 --- structure of
        machine code (x86, but the ideas transfer).
  \item CPUlator (browser RISC-V simulator) --- run every program in
        this chapter interactively.
        \url{https://cpulator.01xz.net/}
  \item Spike (RISC-V ISA simulator) --- the canonical command-line
        reference implementation.
        \url{https://github.com/riscv-software-src/riscv-isa-sim}
\end{itemize}

\section*{Watch instead}
\label{sec:watch03}
\addcontentsline{toc}{section}{Watch instead}

\begin{rulebox}
\textbf{Watch instead.}
\begin{itemize}
  \item ``DDCA Ch6 --- Part 2: RISC-V Instructions'' --- Sarah
        Harris, 2:48. The instructions-first pitch that matches this
        chapter's structure; the book's own DDCA source, on camera.
        \url{https://youtu.be/SdjPQtMPPmo}
  \item ``RISC-V Assembly Hello World (Part 1)'' --- LaurieWired,
        11:42. A first full program assembled and run --- the same
        journey ch03\_sum.S takes on the book's simulator.
        \url{https://youtu.be/0IeOaiKszLk}
\end{itemize}
\end{rulebox}

\chapternotes
